% ============================================================
%  Pareto-Scalarization Variant Scheduling in Serverledge
%  Documento LaTeX compatibile con Overleaf
% ============================================================
\documentclass[a4paper,12pt]{article}

\usepackage[utf8]{inputenc}
\usepackage[T1]{fontenc}
\usepackage[italian]{babel}
\usepackage{geometry}
\geometry{margin=2.5cm}

\usepackage{amsmath, amssymb}
\usepackage{booktabs}
\usepackage{graphicx}
\usepackage{xcolor}
\usepackage{listings}
\usepackage{hyperref}
\usepackage{caption}
\usepackage{float}
\usepackage{enumitem}
\usepackage{tikz}
\usetikzlibrary{arrows.meta, positioning, shapes}

% ── Stile codice ─────────────────────────────────────────────
\definecolor{codebg}{rgb}{0.96,0.96,0.96}
\definecolor{keyword}{rgb}{0.0,0.3,0.7}
\definecolor{comment}{rgb}{0.3,0.5,0.3}
\definecolor{string}{rgb}{0.6,0.1,0.1}

\lstdefinestyle{go}{
  language=C,          % Go non è predefinito; C è sufficiente per evidenziare
  backgroundcolor=\color{codebg},
  basicstyle=\ttfamily\small,
  keywordstyle=\color{keyword}\bfseries,
  commentstyle=\color{comment}\itshape,
  stringstyle=\color{string},
  breaklines=true, breakatwhitespace=true,
  frame=single, framesep=4pt,
  numbers=left, numberstyle=\tiny\color{gray},
  tabsize=4,
  morekeywords={func, type, struct, return, if, for, range, nil,
                var, make, append, len, bool, float64, string, int,
                package, import, error},
  literate={λ}{$\lambda$}{1}
}

\lstdefinestyle{bash}{
  language=bash,
  backgroundcolor=\color{codebg},
  basicstyle=\ttfamily\footnotesize,
  breaklines=true,
  frame=single, framesep=4pt,
  numbers=none,
}

\lstdefinestyle{json}{
  basicstyle=\ttfamily\footnotesize,
  backgroundcolor=\color{codebg},
  breaklines=true,
  frame=single, framesep=4pt,
  numbers=none,
  morestring=[b]",
  stringstyle=\color{string},
}

% ── Intestazione ─────────────────────────────────────────────
\title{\textbf{Pareto-Scalarization Variant Scheduling}\\[6pt]
       \large Progettazione e implementazione in Serverledge}
\author{Relazione tecnica --- Tesi}
\date{Febbraio 2026}

% ============================================================
\begin{document}
\maketitle
\tableofcontents
\newpage

% ============================================================
\section{Motivazione e obiettivi}
% ============================================================

La piattaforma FaaS \emph{Serverledge} supporta il concetto di
\textbf{variante di funzione}: per un unico nome logico (es.\
\texttt{PiLeibniz}) possono esistere più implementazioni registrate,
ognuna con un proprio profilo di energia e di precisione del risultato.

La politica di scheduling originale era basata su due parametri
binari/scalari:
\begin{itemize}
  \item \texttt{AllowApprox} --- abilita la selezione di varianti
        approssimate;
  \item \texttt{MaxEnergyJoule} --- budget energetico massimo (J).
\end{itemize}

Questa politica presentava due limitazioni principali:
\begin{enumerate}
  \item Il parametro \texttt{MaxEnergyJoule} tagliava l'insieme
        ammissibile ma non esprimeva alcuna preferenza relativa tra
        \emph{energia} e \emph{qualità del risultato}.
  \item Non era possibile ottenere un trade-off continuo tra i due
        obiettivi: l'utente non poteva indicare ``preferirei risparmiare
        energia, ma non a qualsiasi costo in termini di precisione''.
\end{enumerate}

\subsection{Obiettivo}

Sostituire l'intera politica con uno schema a
\textbf{scalarizzazione Pareto}, in cui l'utente fornisce un unico
parametro $\lambda \in [0,1]$:
\begin{equation}\label{eq:meaning}
  \lambda = 0 \;\Rightarrow\; \text{minimizza l'energia},
  \qquad
  \lambda = 1 \;\Rightarrow\; \text{minimizza l'errore (massimizza qualità)}.
\end{equation}
Valori intermedi esprimono un trade-off pesato e continuo.

% ============================================================
\section{Algoritmo di selezione Pareto}
% ============================================================

\subsection{Panoramica}

L'algoritmo si articola in cinque passi, implementati in
\texttt{internal/scheduling/variant\_selection.go}.

\begin{enumerate}
  \item \textbf{Raccolta candidati}: vengono caricati da etcd tutti i
        record con lo stesso \texttt{LogicalName} della funzione
        invocata.

  \item \textbf{Valutazione}: per ogni variante si calcolano due
        misure obiettivo:
        \begin{itemize}
          \item $E_i$ --- energia di invocazione in Joule
                (\texttt{invocation\_joule}); si usa sempre
                \emph{warm}, per escludere l'overhead di cold-start
                che è transitorio e dipendente dal runtime, non
                dall'algoritmo \textsection\ref{sec:coldstart};
          \item $\varepsilon_i$ --- stima dell'errore del risultato
                (campo \texttt{error\_estimate} del profilo
                \texttt{OutputModel}).
        \end{itemize}

  \item \textbf{Filtro Pareto}: si rimuovono le varianti
        \emph{dominate}. La variante $j$ domina $i$ se:
        \begin{equation}\label{eq:dom}
          E_j \le E_i \;\land\; \varepsilon_j \le \varepsilon_i
          \;\land\; (E_j < E_i \;\lor\; \varepsilon_j < \varepsilon_i).
        \end{equation}
        Il sottoinsieme delle varianti non dominate costituisce il
        \textbf{fronte di Pareto}.

  \item \textbf{Normalizzazione min-max} sul fronte:
        \begin{align}
          \hat{E}_i &= \frac{E_i - E_{\min}}{E_{\max} - E_{\min}},
          \qquad
          \hat{\varepsilon}_i = \frac{\varepsilon_i - \varepsilon_{\min}}
                                     {\varepsilon_{\max} - \varepsilon_{\min}}.
        \end{align}
        Se il range è zero (tutte le varianti hanno lo stesso valore),
        il termine normalizzato è 0.

  \item \textbf{Scalarizzazione lineare}: per ogni variante del fronte:
        \begin{equation}\label{eq:score}
          \text{score}(i) = (1-\lambda)\,\hat{E}_i + \lambda\,\hat{\varepsilon}_i.
        \end{equation}
        Si restituisce la variante con \textbf{score minimo}.
\end{enumerate}

\subsection{Energia effettiva: cold-start condizionale}\label{sec:coldstart}

Il costo di cold-start ($C_{\text{cs}}$) è l'energia necessaria per
avviare un nuovo container quando non è disponibile un'istanza calda.
A differenza di quanto ipotizzato inizialmente, questo costo
\emph{non viene ignorato} nella comparazione Pareto: il criterio
adottato è che la valutazione rifletta il \textbf{costo reale}
dell'invocazione imminente.

\begin{equation}\label{eq:effective}
  E_i^{\text{eff}} =
  \begin{cases}
    E_i^{\text{inv}} & \text{se esiste un container caldo per la variante } i \\
    C_{\text{cs},i} + E_i^{\text{inv}} & \text{altrimenti (cold start)}
  \end{cases}
\end{equation}

Il fronte di Pareto è quindi \emph{state-dependent}: dipende da quali
container sono in esecuzione al momento dell'invocazione.  Una variante
che risulterebbe dominata se fosse cold può diventare Pareto-ottimale
una volta che il suo container è già caldo, e viceversa.

\paragraph{Effetto pratico.}
Nel caso di PiLeibniz, le varianti Python in cold-start pagano
$C_{\text{cs,Py}} \approx 0.069{-}0.071\,\text{J}$, che domina di
gran lunga il costo di invocazione ($\sim 0.0002{-}0.008\,\text{J}$).
Al contrario, la variante C ha $C_{\text{cs,C}} \approx 0.009\,\text{J}$.
Il fronte di Pareto \emph{a freddo} può quindi differire
sostanzialmente da quello \emph{a caldo}, incentivando il riuso
dei container come meccanismo di risparmio energetico.

\subsection{Algoritmo di filtro: sort + scan in $O(n \log n)$}

L'implementazione adotta un algoritmo ottimale per il caso a 2 obiettivi,
che riduce la complessità da $O(n^2)$ (doppio ciclo naive) a
$O(n \log n)$.

L'intuizione chiave è che la relazione di dominanza, su due assi,
può essere sfruttata dopo un ordinamento: se l'insieme è già
ordinate per energia crescente, un punto può essere dominato
solo da un punto che lo precede nell'ordinamento (che ha quindi
energia $\le$). Questo riduce il controllo di dominanza a
un semplice confronto sull'asse residuo.

\paragraph{Passo 1 --- Ordinamento.}
I candidati vengono ordinati per $E$ crescente.
A parità di $E$, si ordina per $\varepsilon$ crescente: tra due
varianti con identica energia, quella con errore maggiore è
automaticamente dominata dall'altra, quindi viene processata
\emph{dopo} e sarà scartata al passo successivo.

\paragraph{Passo 2 --- Scansione singola.}
Si percorre la lista ordinata da sinistra a destra mantenendo una
variabile $\varepsilon_{\min}$, inizializzata a $+\infty$, che
registra il \_minimo errore visto finora sul fronte.

Un punto $p$ è Pareto-ottimale se e solo se
$\varepsilon_p < \varepsilon_{\min}$.

\textit{Perché è corretto}: tutti i punti già processati hanno
$E \le E_p$ (per la proprietà dell'ordinamento).
Se esistesse un punto precedente $q$ che domina $p$, si avrebbe
$E_q \le E_p$ e $\varepsilon_q \le \varepsilon_p$, con almeno uno
stretto. Poiché $E_q \le E_p$ è già garantito, la condizione di
dominanza si riduce a $\varepsilon_q \le \varepsilon_p$.
L'unica configurazione in cui $p$ \emph{non} è dominato è quindi
$\varepsilon_p < \varepsilon_{\min}$: nessun punto precedente ha
errore $\le \varepsilon_p$.

\paragraph{Esempio su PiLeibniz.}
Dopo ordinamento per $E$ crescente, la scansione procede:

\begin{table}[H]
\centering
\small
\begin{tabular}{@{}lrrrrc@{}}
\toprule
Variante & $E$ (J) & $\varepsilon$ & $\varepsilon_{\min}$ pre & $\varepsilon_p < \varepsilon_{\min}$? & Fronte \\
\midrule
n1000   & 0.000200 & 0.001000  & $+\infty$ & \checkmark & \checkmark \\
n5000   & 0.000300 & 0.000200  & 0.001000  & \checkmark & \checkmark \\
n10000  & 0.000430 & 0.000100  & 0.000200  & \checkmark & \checkmark \\
n50000  & 0.000800 & 0.000015  & 0.000100  & \checkmark & \checkmark \\
c       & 0.001037 & 0.000000  & 0.000015  & \checkmark & \checkmark \\
n100000 & 0.001150 & 0.000008  & 0.000000  & $\times$   & --- \\
n200000 & 0.001550 & 0.000003  & 0.000000  & $\times$   & --- \\
opt-py  & 0.007400 & 0.000000  & 0.000000  & $\times$   & --- \\
base    & 0.007750 & 0.000000  & 0.000000  & $\times$   & --- \\
\bottomrule
\end{tabular}
\caption{Traccia dell'algoritmo sort+scan su PiLeibniz. Le ultime 4 righe
         vengono scartate perché $\varepsilon_{\min}$ è già 0 dopo \texttt{c}.
         Il risultato è identico all'algoritmo naive $O(n^2)$.}
\end{table}

\begin{lstlisting}[style=go, caption={Scalarizzazione --- \texttt{SelectParetoVariant()}}]
for _, v := range pareto {
    var normE, normErr float64
    if eRange > 0   { normE   = (v.energy   - eMin) / eRange }
    if errRange > 0 { normErr = (v.errScore - errMin) / errRange }

    score := (1-lambda)*normE + lambda*normErr
    if score < bestScore {
        bestScore = score
        bestV = v
    }
}
\end{lstlisting}

% ============================================================
\section{Modifiche al codice sorgente}
% ============================================================

\subsection{File modificati}

\begin{table}[H]
\centering
\caption{Riepilogo dei file modificati durante l'implementazione.}
\begin{tabular}{@{}lp{9cm}@{}}
\toprule
\textbf{File} & \textbf{Modifica principale} \\
\midrule
\texttt{internal/function/request.go}
  & Rimossi \texttt{AllowApprox}, \texttt{MaxEnergyJoule};
    aggiunto \texttt{QualityWeight *float64},
    struct \texttt{ParetoPoint}, campo \texttt{ParetoFront} in
    \texttt{VariantSchedulingReport}. \\
\addlinespace
\texttt{internal/client/types.go}
  & \texttt{InvocationRequest.QualityWeight *float64} (tag JSON
    \texttt{qualityWeight}); rimossi i vecchi campi. \\
\addlinespace
\texttt{internal/cli/cli.go}
  & Flag \texttt{--quality-weight float64} (default $-1$);
    rimossi \texttt{--allowApprox} e \texttt{--maxEnergyJoule}. \\
\addlinespace
\texttt{internal/api/api.go}
  & Passaggio di \texttt{QualityWeight} dall'oggetto HTTP alla
    \texttt{Request}; condizione di trigger \texttt{!= nil}. \\
\addlinespace
\texttt{internal/scheduling/scheduler.go}
  & Condizione di trigger cambiata in
    \texttt{r.QualityWeight != nil};
    chiamata a \texttt{SelectParetoVariant}. \\
\addlinespace
\texttt{internal/scheduling/variant\_selection.go}
  & Riscrittura completa: \texttt{SelectParetoVariant()} +
    \texttt{paretoFilter()}. \\
\addlinespace
\texttt{variants/PiLeibniz/PiLeibniz.json}
  & 9 varianti (5 Pareto-ottimali + 4 dominate per test). \\
\addlinespace
\texttt{variants/fibonacci/fibonacci.json}
  & Invariato; 3 varianti, tutte con \texttt{error\_estimate=0.0}. \\
\bottomrule
\end{tabular}
\end{table}

\subsection{Strutture dati principali}

\begin{lstlisting}[style=go, caption={Nuove strutture in \texttt{request.go}}]
// Aggiunto al tipo Request:
QualityWeight *float64  // nil = nessuna selezione variante

// Nuovo punto sul fronte di Pareto:
type ParetoPoint struct {
    FunctionName  string  `json:"function_name"`
    VariantID     string  `json:"variant_id,omitempty"`
    Energy        float64 `json:"energy_joule"`
    ErrorEstimate float64 `json:"error_estimate"`
    NormEnergy    float64 `json:"norm_energy"`
    NormError     float64 `json:"norm_error"`
    Score         float64 `json:"score"`
}

// Report restituito nella risposta HTTP:
type VariantSchedulingReport struct {
    LogicalName     string       `json:"logical_name"`
    InvokedFunction string       `json:"invoked_function"`
    SelectedFunction string      `json:"selected_function"`
    VariantID       string       `json:"variant_id,omitempty"`
    QualityWeight   float64      `json:"quality_weight"`
    EstimatedEnergy float64      `json:"estimated_energy_joule,omitempty"`
    WarmHint        bool         `json:"warm_hint"`
    ErrorEstimate   float64      `json:"error_estimate"`
    DecisionReason  string       `json:"decision_reason,omitempty"`
    ParetoFront     []ParetoPoint `json:"pareto_front,omitempty"`
}
\end{lstlisting}

% ============================================================
\section{Interfaccia utente}
% ============================================================

\subsection{CLI}

\begin{lstlisting}[style=bash, caption={Invocazione tramite CLI}]
# Minimizza energia (lambda = 0)
./serverledge-cli invoke --function PiLeibniz --quality-weight 0.0

# Trade-off bilanciato (lambda = 0.5)
./serverledge-cli invoke --function PiLeibniz --quality-weight 0.5

# Minimizza errore / massimizza qualita' (lambda = 1)
./serverledge-cli invoke --function PiLeibniz --quality-weight 1.0

# Senza quality-weight: scheduling tradizionale (nessuna selezione variante)
./serverledge-cli invoke --function PiLeibniz --param n:1000
\end{lstlisting}

\subsection{API HTTP (JSON)}

\begin{lstlisting}[style=json, caption={Corpo della richiesta HTTP}]
POST /invoke/PiLeibniz
{
  "qualityWeight": 0.5,
  "params": {}
}
\end{lstlisting}

\subsection{Formato della risposta}

La risposta include un campo \texttt{variant\_scheduling} con il fronte
di Pareto completo, utile per produrre grafici energia/errore:

\begin{lstlisting}[style=json, caption={Estratto della risposta JSON}]
{
  "variant_scheduling": {
    "logical_name": "PiLeibniz",
    "selected_function": "PiLeibniz-n5000",
    "variant_id": "n5000",
    "quality_weight": 0.5,
    "estimated_energy_joule": 0.0003,
    "error_estimate": 0.0002,
    "decision_reason": "pareto-scalarisation",
    "pareto_front": [
      { "function_name": "PiLeibniz-c",     "energy_joule": 0.001037, "error_estimate": 0.0,    "norm_energy": 1.0,    "norm_error": 0.0,   "score": 0.5 },
      { "function_name": "PiLeibniz-n1000", "energy_joule": 0.0002,   "error_estimate": 0.001,  "norm_energy": 0.0,    "norm_error": 1.0,   "score": 0.5 },
      { "function_name": "PiLeibniz-n5000", "energy_joule": 0.0003,   "error_estimate": 0.0002, "norm_energy": 0.1195, "norm_error": 0.2,   "score": 0.1597 },
      { "function_name": "PiLeibniz-n10000","energy_joule": 0.00043,  "error_estimate": 0.0001, "norm_energy": 0.2748, "norm_error": 0.1,   "score": 0.1874 },
      { "function_name": "PiLeibniz-n50000","energy_joule": 0.0008,   "error_estimate": 0.000015,"norm_energy": 0.7168,"norm_error": 0.015, "score": 0.3659 }
    ]
  }
}
\end{lstlisting}

% ============================================================
\section{Caso studio: PiLeibniz}
% ============================================================

\subsection{Varianti registrate}

Il calcolo di $\pi$ con la serie di Leibniz:
\begin{equation}
  \pi = 4\sum_{k=0}^{n-1} \frac{(-1)^k}{2k+1}
\end{equation}
è approssimato a piacere variando $n$.  Sono state definite 9 varianti,
riassunte nella Tabella~\ref{tab:varianti}.

\paragraph{Convenzione sull'errore per le varianti baseline.}
Le varianti \texttt{base}, \texttt{optimization-py} e \texttt{c} usano
$n = 1\,000\,000$ iterazioni fisse, producendo un errore teorico
$\varepsilon \approx 1/n = 10^{-6}$.
Nel JSON il campo \texttt{error\_estimate} è impostato a \texttt{0.0} per
convenzione: questi implementazioni sono considerate la
\emph{baseline di riferimento} per la qualità, e l'errore delle varianti
approssimate è espresso \emph{relativamente} ad esse.
Questo non influenza l'algoritmo Pareto (che legge i valori direttamente
dal JSON), ma va tenuto presente nell'interpretazione dei risultati.

\paragraph{Cold start per varianti dello stesso runtime.}
Tutte le varianti \texttt{python310} girano nello stesso tipo di container
Docker, quindi il loro \texttt{cold\_start\_joule} deve essere identico:
$C_{\text{cs, py}} = 0.069905$\,J per tutte.
Il valore è diverso solo per il runtime \texttt{native} (C/Alpine):
$C_{\text{cs, c}} = 0.008653$\,J, giustificato dal container più leggero.
Un \texttt{cold\_start\_joule} difforte tra varianti dello stesso runtime
indica dati di profilo incoerenti e, seppur non influenzi la selezione
Pareto (che usa esclusivamente \texttt{invocation\_joule}), invaliderebbe
il calcolo del costo totale nel path di cold-start dello scheduler.

\begin{table}[H]
\centering
\caption{Varianti di PiLeibniz ($E$ = \texttt{invocation\_joule}).}
\label{tab:varianti}
\begin{tabular}{@{}llrrcc@{}}
\toprule
\textbf{ID} & \textbf{Runtime} & $E$ (J) & $\varepsilon$ & Pareto? & Note \\
\midrule
n1000    & python310 & 0.000200 & 0.001000 & \checkmark & \\
n5000    & python310 & 0.000300 & 0.000200 & \checkmark & \\
n10000   & python310 & 0.000430 & 0.000100 & \checkmark & \\
n50000   & python310 & 0.000800 & 0.000015 & \checkmark & \\
c        & native    & 0.001037 & 0.000000 & \checkmark & \\
\midrule
n100000  & python310 & 0.001150 & 0.000008 & $\times$ & dominata da \texttt{c} \\
n200000  & python310 & 0.001550 & 0.000003 & $\times$ & dominata da \texttt{c} \\
base     & python310 & 0.007750 & 0.000000 & $\times$ & dominata da \texttt{c} \\
opt-py   & python310 & 0.007400 & 0.000000 & $\times$ & dominata da \texttt{c} \\
\bottomrule
\end{tabular}
\end{table}

\subsubsection{Analisi delle dominanze}

\begin{itemize}
  \item \texttt{n100000}: $E=0.001150 > E_c=0.001037$ e
        $\varepsilon=0.000008 > \varepsilon_c=0$ $\Rightarrow$ dominata da \texttt{c}.
  \item \texttt{n200000}: $E=0.001550 > E_c$ e
        $\varepsilon=0.000003 > \varepsilon_c$ $\Rightarrow$ dominata da \texttt{c}.
  \item \texttt{base}: $E=0.007750 \gg E_c$ e $\varepsilon=0$ $\Rightarrow$
        dominata da \texttt{c} (stessa qualità, molto più costosa).
  \item \texttt{opt-py}: ancora più costosa di \texttt{base}, dominata da
        \texttt{c}.
\end{itemize}

Queste 4 varianti sono mantenute nel JSON intenzionalmente come
\emph{test data} per verificare che il filtro Pareto le escluda
correttamente.

\subsection{Fronte di Pareto e selezione per vari $\lambda$}

\begin{figure}[H]
\centering
\begin{tikzpicture}[scale=1]
  % Assi
  \draw[->] (0,0) -- (7,0) node[right] {Energia $E$ (J $\times 10^{-4}$)};
  \draw[->] (0,0) -- (0,6) node[above] {Errore $\varepsilon$};

  % Pareto front (punti normalizzati per il disegno)
  \coordinate (n1000)  at (1.0, 5.0); % E=0.2, err=1.000
  \coordinate (n5000)  at (1.5, 2.0); % E=0.3, err=0.200
  \coordinate (n10000) at (2.15,1.0); % E=0.43,err=0.100
  \coordinate (n50000) at (4.0, 0.15);% E=0.8, err=0.015
  \coordinate (c)      at (5.185,0.0);% E=1.037,err=0.0

  % Punti dominati (grigi)
  \filldraw[gray!50] (5.75, 0.08) circle (3pt) node[right,gray,small] {\small n100000};
  \filldraw[gray!50] (7.75, 0.03) circle (3pt);
  \filldraw[gray!50] (38.75,0.0)  circle (0pt);% troppo lontano, non disegnare

  % Linea fronte
  \draw[blue, thick] (n1000) -- (n5000) -- (n10000) -- (n50000) -- (c);

  % Punti Pareto
  \foreach \p/\lbl in {n1000/n1000, n5000/n5000, n10000/n10000, n50000/n50000, c/c} {
    \filldraw[blue] (\p) circle (3pt);
    \node[above right, blue, font=\small] at (\p) {\texttt{\lbl}};
  }

  % Annotazioni lambda
  \draw[dashed,red] (1.0,5.0) -- ++(0,-0.5) node[below,red,font=\scriptsize]{$\lambda{=}0$};
  \draw[dashed,orange] (1.5,2.0) -- ++(0.3,0.5) node[above,orange,font=\scriptsize]{$\lambda{=}0.2{,}0.5$};
  \draw[dashed,darkgreen] (5.185,0.0) -- ++(-0.3,0.5) node[above,darkgreen,font=\scriptsize]{$\lambda{=}1$};
\end{tikzpicture}
\caption{Fronte di Pareto per PiLeibniz (5 punti non dominati).
         I quadranti mostrano la direzione di dominanza: un punto è ottimale se
         nessun altro si trova in basso a sinistra rispetto ad esso.}
\end{figure}

\subsubsection{Calcolo degli score per $\lambda = 0.5$}

Con $E_{\min}=0.0002$, $E_{\max}=0.001037$, $\varepsilon_{\min}=0$,
$\varepsilon_{\max}=0.001$:

\begin{table}[H]
\centering
\begin{tabular}{@{}lrrrr@{}}
\toprule
Variante & $\hat{E}$ & $\hat{\varepsilon}$ & $\text{score}_{0.5}$ & \\
\midrule
n1000   & 0.000 & 1.000 & 0.500 & \\
n5000   & 0.119 & 0.200 & \textbf{0.160} & $\leftarrow$ min \\
n10000  & 0.275 & 0.100 & 0.187 & \\
n50000  & 0.717 & 0.015 & 0.366 & \\
c       & 1.000 & 0.000 & 0.500 & \\
\bottomrule
\end{tabular}
\caption{Score per $\lambda=0.5$; vince \texttt{n5000}.}
\end{table}

\subsection{Risultati dei test --- output di terminale}

I seguenti output sono stati ottenuti eseguendo
\texttt{ScriptTesi/PiLeibniz/SchedulingPiLiebniz.sh} con il server
Serverledge attivo.

\subsubsection{Invocazione 1 --- $\lambda = 0.0$ (min energia)}

Risultato atteso: \texttt{PiLeibniz-n1000} (punto più a sinistra sul fronte).

\begin{lstlisting}[style=json]
"variant_scheduling": {
    "logical_name": "PiLeibniz",
    "selected_function": "PiLeibniz-n1000",   ← CORRETTO
    "variant_id": "n1000",
    "quality_weight": 0,
    "estimated_energy_joule": 0.0002,
    "error_estimate": 0.001,
    "decision_reason": "pareto-scalarisation",
    "pareto_front": [
        {"function_name":"PiLeibniz-c",     "norm_energy":1, "norm_error":0, "score":1},
        {"function_name":"PiLeibniz-n1000", "norm_energy":0, "norm_error":1, "score":0},
        {"function_name":"PiLeibniz-n5000", "norm_energy":0.1195,"norm_error":0.2,"score":0.1195},
        {"function_name":"PiLeibniz-n10000","norm_energy":0.2748,"norm_error":0.1,"score":0.2748},
        {"function_name":"PiLeibniz-n50000","norm_energy":0.7168,"norm_error":0.015,"score":0.7168}
    ]
}
\end{lstlisting}

\textit{Spiegazione}: con $\lambda=0$ lo score si riduce a
$\text{score}(i) = \hat{E}_i$.  La variante con energia minima
(\texttt{n1000}, $E=0.0002$\,J) ha $\hat{E}=0$, punteggio~0.

\subsubsection{Invocazione 2 --- $\lambda = 1.0$ (min errore)}

Risultato atteso: \texttt{PiLeibniz-c} (unico punto con $\varepsilon=0$).

\begin{lstlisting}[style=json]
"variant_scheduling": {
    "selected_function": "PiLeibniz-c",   ← CORRETTO
    "variant_id": "c",
    "quality_weight": 1,
    "estimated_energy_joule": 0.001037,
    "error_estimate": 0,
    "pareto_front": [
        {"function_name":"PiLeibniz-c",     "norm_energy":1,"norm_error":0,"score":0},
        {"function_name":"PiLeibniz-n1000", "norm_energy":0,"norm_error":1,"score":1},
        ...
    ]
}
\end{lstlisting}

\textit{Spiegazione}: con $\lambda=1$ lo score è $\hat{\varepsilon}_i$.
Solo \texttt{c} ha errore~0 (esatto per natura dell'algoritmo
in virgola mobile), quindi punteggio~0.

\subsubsection{Invocazione 3 --- $\lambda = 0.5$ (trade-off bilanciato)}

Risultato atteso: \texttt{PiLeibniz-n5000}.

\begin{lstlisting}[style=json]
"variant_scheduling": {
    "selected_function": "PiLeibniz-n5000",   ← CORRETTO
    "quality_weight": 0.5,
    "estimated_energy_joule": 0.0003,
    "error_estimate": 0.0002,
    "pareto_front": [
        {"function_name":"PiLeibniz-n5000","norm_energy":0.1195,"norm_error":0.2,"score":0.1597}
        ...
    ]
}
\end{lstlisting}

\textit{Spiegazione}: \texttt{n5000} minimizza la combinazione
$(0.5 \cdot \hat{E} + 0.5 \cdot \hat{\varepsilon})$ grazie a un
buon compromesso tra energia (solo $+50\%$ rispetto al minimo) e
errore (solo $1/5$ di quello di \texttt{n1000}).

\subsubsection{Invocazione 4 --- $\lambda = 0.2$ (priorità all'energia)}

Risultato atteso: \texttt{PiLeibniz-n5000}.

\begin{lstlisting}[style=json]
"variant_scheduling": {
    "selected_function": "PiLeibniz-n5000",   ← CORRETTO
    "quality_weight": 0.2,
    "estimated_energy_joule": 0.0003,
    "error_estimate": 0.0002
}
\end{lstlisting}

\textit{Spiegazione}: anche con il peso dell'errore ridotto a 0.2,
\texttt{n5000} rimane il punto con score minimo (0.136). Sebbene
\texttt{n1000} abbia $\hat{E}=0$, il suo errore alto
($\hat{\varepsilon}=1$) contribuisce $0.2 \times 1 = 0.2$, superiore
allo score 0.136 di \texttt{n5000}.

\subsubsection{Riepilogo PiLeibniz}

\begin{table}[H]
\centering
\caption{Risultati delle 4 invocazioni PiLeibniz.}
\begin{tabular}{@{}crll@{}}
\toprule
$\lambda$ & Atteso & Selezionato & Corretto? \\
\midrule
0.0 & n1000  & \texttt{PiLeibniz-n1000} & \checkmark \\
1.0 & c      & \texttt{PiLeibniz-c}     & \checkmark \\
0.5 & n5000  & \texttt{PiLeibniz-n5000} & \checkmark \\
0.2 & n5000  & \texttt{PiLeibniz-n5000} & \checkmark \\
\bottomrule
\end{tabular}
\end{table}

% ============================================================
\section{Caso studio: Fibonacci (caso degenere)}
% ============================================================

\subsection{Varianti registrate}

\begin{table}[H]
\centering
\caption{Varianti di Fibonacci.}
\begin{tabular}{@{}llrrcc@{}}
\toprule
\textbf{ID} & \textbf{Runtime} & $E$ (J) & $\varepsilon$ & Pareto? \\
\midrule
base         & python310 & 0.000732 & 0.0 & $\times$ \\
opt-py       & python310 & 0.000495 & 0.0 & $\times$ \\
c            & native    & 0.000172 & 0.0 & \checkmark \\
\bottomrule
\end{tabular}
\end{table}

\subsection{Analisi del fronte di Pareto}

Tutte e tre le varianti hanno $\varepsilon = 0$ (Fibonacci è un
calcolo esatto).  Con errore identico, il confronto si riduce
all'asse energetico:
\begin{equation}
  E_c = 0.000172 < E_{\text{opt-py}} = 0.000495 < E_{\text{base}} = 0.000732.
\end{equation}
La variante \texttt{c} domina entrambe le Python (Eq.~\ref{eq:dom}),
poiché ha energia strettamente inferiore e uguale errore.

Il \textbf{fronte di Pareto è composto da un singolo punto}:
$\mathcal{P} = \{\texttt{c}\}$.

\subsection{Caso degenere: fronte con un solo punto}

Quando il fronte ha un solo elemento, la normalizzazione min-max
produce $\hat{E}=0$ e $\hat{\varepsilon}=0$, e lo
score è 0 per qualsiasi $\lambda$. L'algoritmo è robusto a questo
caso: non è richiesta nessuna gestione speciale.

\begin{lstlisting}[style=go, caption={Gestione robusta quando range = 0}]
if eRange > 0   { normE   = (v.energy   - eMin) / eRange } // else 0
if errRange > 0 { normErr = (v.errScore - errMin) / errRange } // else 0
\end{lstlisting}

\subsection{Risultati dei test --- output di terminale}

I seguenti output sono stati ottenuti eseguendo
\texttt{ScriptTesi/Fibonacci/SchedulerFibonacci.sh}.

\subsubsection{Invocazione 1 --- $\lambda = 0.0$}

\begin{lstlisting}[style=json]
"variant_scheduling": {
    "logical_name": "fibonacci",
    "selected_function": "fibonacci-c",   ← unico punto Pareto
    "variant_id": "c",
    "quality_weight": 0,
    "estimated_energy_joule": 0.000171579,
    "error_estimate": 0,
    "decision_reason": "pareto-scalarisation",
    "pareto_front": [
        {
            "function_name": "fibonacci-c",
            "energy_joule": 0.000171579,
            "error_estimate": 0,
            "norm_energy": 0,
            "norm_error": 0,
            "score": 0
        }
    ]
}
\end{lstlisting}

\subsubsection{Invocazione 2 --- $\lambda = 0.5$ (warm start)}

\begin{lstlisting}[style=json]
"variant_scheduling": {
    "selected_function": "fibonacci-c",   ← invariato
    "warm_hint": true,
    "quality_weight": 0.5,
    "pareto_front": [
        {"function_name":"fibonacci-c","norm_energy":0,"norm_error":0,"score":0}
    ]
}
\end{lstlisting}

\subsubsection{Invocazione 3 --- $\lambda = 1.0$}

\begin{lstlisting}[style=json]
"variant_scheduling": {
    "selected_function": "fibonacci-c",   ← invariato
    "warm_hint": true,
    "quality_weight": 1,
    "pareto_front": [
        {"function_name":"fibonacci-c","norm_energy":0,"norm_error":0,"score":0}
    ]
}
\end{lstlisting}

\subsubsection{Riepilogo Fibonacci}

\begin{table}[H]
\centering
\caption{Risultati delle 3 invocazioni Fibonacci.}
\begin{tabular}{@{}crll@{}}
\toprule
$\lambda$ & Atteso & Selezionato & Corretto? \\
\midrule
0.0 & c & \texttt{fibonacci-c} & \checkmark \\
0.5 & c & \texttt{fibonacci-c} & \checkmark \\
1.0 & c & \texttt{fibonacci-c} & \checkmark \\
\bottomrule
\end{tabular}
\end{table}

Il fronte di Pareto con un singolo punto è il
\textbf{caso degenere} della scalarizzazione lineare: indipendentemente
da $\lambda$, la scelta è univoca e corretta.

% ============================================================
\section{Nota tecnica: compilazione binari nativi}
% ============================================================

Il runtime \texttt{native} di Serverledge esegue i binari dentro
container Alpine Linux, che usa la libreria C musl anziché glibc.
I binari compilati con le toolchain standard su Ubuntu/Debian sono
linked dinamicamente contro glibc e falliscono con
\texttt{not found} dentro Alpine.

La soluzione è la \textbf{compilazione statica}:

\begin{lstlisting}[style=bash]
# Fibonacci (gia' presente staticamente linkato)
gcc -O2 -static -o fibonacci fibonacci.c

# PiLeibniz C (compilato durante lo sviluppo)
gcc -O2 -static -o piLeibniz_base piLeibniz_base.c
\end{lstlisting}

Il binario risultante include musl tramite static linking e gira
in qualsiasi container Linux senza dipendenze esterne.

% ============================================================
\section{Conclusioni}
% ============================================================

La politica di scheduling basata su scalarizzazione Pareto introduce
un controllo continuo e principiato del trade-off energetico/qualitativo
nelle invocazioni FaaS:

\begin{enumerate}
  \item \textbf{Convergenza ottimale}: la scalarizzazione lineare
        copre l'intero fronte di Pareto al variare di $\lambda$.
  \item \textbf{Robustezza}: il caso degenere (fronte con un punto solo)
        è gestito senza eccezioni.
  \item \textbf{Trasparenza}: il campo \texttt{pareto\_front} nella
        risposta permette all'utente di ispezionare tutti i punti
        ottimali e i rispettivi score.
  \item \textbf{Esclusione del cold-start}: la comparazione utilizza
        solo l'energia di invocazione (warm), evitando distorsioni
        dipendenti dal runtime.
\end{enumerate}

I test su due funzioni reali (PiLeibniz con 9 varianti e Fibonacci con
3) confermano il comportamento corretto per tutti i valori di $\lambda$
testati.

\end{document}
